\documentclass[11pt, a4paper]{article}

\usepackage[utf8]{inputenc}
\usepackage[T1]{fontenc}
\usepackage{amsmath}
\usepackage{amssymb}
\usepackage{graphicx}
\usepackage{geometry}
\usepackage{fancyhdr}
\usepackage{enumitem}
\usepackage{array}
\usepackage{eso-pic} % Required for placing the logo
\usepackage{esvect}
\usepackage[most]{tcolorbox} % Needed for solution box

\usepackage{siunitx}
\DeclareSIUnit{\litre}{l}
\sisetup{
    locale = DE,
    inter-unit-product = \cdot,
    per-mode = symbol-or-fraction
}

\makeatletter
\def\input@path{{./}{../../}}
\makeatother

\geometry{a4paper, top=2.5cm, bottom=2.5cm, left=2.5cm, right=2.5cm}

\input{defs}

% --- Lösungsbox-Umgebung ---
\definecolor{boxcol_back_lgreen}{RGB}{220, 255, 220}
\definecolor{boxcol_back_white}{RGB}{255, 255, 255}
\definecolor{boxcol_frame_green}{RGB}{30, 180, 30}
\newtcolorbox{solutionbox}[1][]{
    colback=boxcol_back_white,
  colframe=boxcol_frame_green,
  colbacktitle=boxcol_back_lgreen,
  fonttitle=\bfseries,
  coltitle=black,
  boxsep=5pt,
  arc=1mm,
  enhanced,
  attach boxed title to top left={yshift=-2mm, xshift=3mm},
  title=Lösung,
  #1
}
% ---

\pagestyle{fancy}
\fancyhf{}
\cfoot{\thepage}

\begin{document}

\section*{\LARGE Lösung: Kurztest 2}
\noindent
\begin{tabular}{ll}
    \textbf{Gruppe: } & \LARGE\textbf{A}  \\[0.3cm]
\end{tabular}

\begin{enumerate}[label=\textbf{\arabic*.},itemsep=0.50cm]

    \item \textbf{Begriffe und Einheiten:} \\
    Nenne die SI-Einheiten von Weg, Zeit, Geschwindigkeit und Beschleunigung.
    \begin{solutionbox}
        \begin{itemize}[itemsep=0pt, topsep=0pt]
            \item Weg: \si{\meter}
            \item Zeit: \si{\second}
            \item Geschwindigkeit: \si{\meter\per\second}
            \item Beschleunigung: \si{\meter\per\second\squared}
        \end{itemize}
    \end{solutionbox}

    \item \textbf{Einheitenumrechnung:} \\
    Ein Fahrrad fährt mit einer Geschwindigkeit von \SI{54}{\kilo\meter\per\hour}. Rechne diese Geschwindigkeit in \unit{\meter\per\second} um und zeige deine Rechnung.
    \begin{solutionbox}
        \[
        \SI{54}{\kilo\meter\per\hour}
        = 54 \cdot \frac{1000}{3600} \si{\meter\per\second}
        = \SI{15}{\meter\per\second}
        \]
        Oder direkt durch $3.6$ teilen: \[54 / 3.6 = \SI{15}{\meter\per\second}.\]
    \end{solutionbox}

    \item \textbf{Gleichförmige Bewegung:} \\
    Ein Auto bewegt sich gleichförmig mit \SI{15}{\meter\per\second}. Welche Strecke legt es in \SI{8}{\second} zurück?
    \begin{solutionbox}
        \[
        s = v \cdot t = \SI{15}{\meter\per\second} \cdot \SI{8}{\second} = \SI{120}{\meter}
        \]
    \end{solutionbox}

    \item \textbf{Durchschnittsgeschwindigkeit:} \\
    Ein Läufer legt \SI{200}{\meter} in \SI{25}{\second} zurück. Berechne seine Durchschnittsgeschwindigkeit.
    \begin{solutionbox}
        \[
        \overline{v} = \frac{\Delta s}{\Delta t} = \frac{\SI{200}{\meter}}{\SI{25}{\second}} = \SI{8}{\meter\per\second}
        \]
    \end{solutionbox}

    \item \textbf{Weg-Zeit-Diagramm:} \\
    Was bedeutet die Steigung in einem Weg-Zeit-Diagramm?
    \begin{solutionbox}
        Die Steigung entspricht der Geschwindigkeit. Eine größere Steigung bedeutet eine größere Geschwindigkeit.
    \end{solutionbox}

    \item \textbf{Beschleunigung:} \\
    Ein Roller startet aus dem Stand und erreicht in \SI{6}{\second} die Geschwindigkeit \SI{12}{\meter\per\second}. Berechne die mittlere Beschleunigung.
    \begin{solutionbox}
        \[
        a = \frac{\Delta v}{\Delta t} = \frac{\SI{12}{\meter\per\second} - 0}{\SI{6}{\second}} = \SI{2}{\meter\per\second\squared}
        \]
    \end{solutionbox}

    \item \textbf{Geradlinige Bewegung mit Startwert:} \\
    Ein Körper befindet sich zur Zeit $t = 0$ an der Position $s_0 = \SI{3}{\meter}$ und bewegt sich mit konstanter Geschwindigkeit \SI{4}{\meter\per\second}. Wo befindet sich der Körper nach \SI{5}{\second}?
    \begin{solutionbox}
        \[
        s(t) = s_0 + v \cdot t = \SI{3}{\meter} + \SI{4}{\meter\per\second} \cdot \SI{5}{\second} = \SI{23}{\meter}
        \]
    \end{solutionbox}

    \item \textbf{Freier Fall:} \\
    Die allgemeine Wegfunktion einer gleichmäßig beschleunigten Bewegung lautet
    \[ s(t) = s_0 + v_0 t + \frac{1}{2} a t^2. \]
    Ein Körper fällt aus \SI{20}{\meter} Höhe. Vernachlässige den Luftwiderstand und verwende $g \approx \SI{10}{\meter\per\second^2}$. Wie lange dauert der freie Fall ungefähr?
    \begin{solutionbox}
        Für den freien Fall gilt $s_0 = \SI{20}{\meter}$, $v_0 = 0$ und $a = -g$.
        Mit der Wegfunktion folgt für den Boden $s(t)=0$:
        \[
        0 = \SI{20}{\meter} - \frac{1}{2} \cdot \SI{10}{\meter\per\second^2} \cdot t^2
        \]
        also
        \[
        t = \sqrt{\frac{2h}{g}} = \sqrt{\frac{2 \cdot \SI{20}{\meter}}{\SI{10}{\meter\per\second^2}}} = \sqrt{4} = \SI{2}{\second}
        \]
    \end{solutionbox}

    \item \textbf{Kreisbewegung:} \\
    Ein Massepunkt bewegt sich auf einem Kreis mit dem Radius \SI{0.5}{\meter}. Für eine volle Umdrehung braucht er \SI{2}{\second}. Welche Winkelgeschwindigkeit hat er und wie groß ist seine Bahngeschwindigkeit?
    \begin{solutionbox}
        Der Umlaufzeitraum ist $T = \SI{2}{\second}$.
        Damit gilt für die Winkelgeschwindigkeit
        \[
        \omega = \frac{2\pi}{T} = \frac{2\pi}{\SI{2}{\second}} = \pi\,\si{\radian\per\second}
        \]
        und für die Bahngeschwindigkeit
        \[
        v = R \cdot \omega = \SI{0.5}{\meter} \cdot \pi\,\si{\radian\per\second} = \frac{\pi}{2}\,\si{\meter\per\second} \approx \SI{1.57}{\meter\per\second}
        \]
    \end{solutionbox}

    \item \textbf{Verständnisfrage:} \\
    In welche Richtung wirkt die Zentripetalbeschleunigung immer? Welche Beschleunigung muss auf einen Massepunkt wirken, wenn er mit der Geschwindigkeit $v = \SI{2}{\meter\per\second}$ im Abstand $r = \SI{0.5}{\meter}$ vom Mittelpunkt kreist?
    \begin{solutionbox}
        Die Zentripetalbeschleunigung wirkt immer zum Mittelpunkt der Kreisbahn.

        Für einen Massepunkt mit Geschwindigkeit $v$ im Abstand $r$ gilt
        \[
        a_{\mathrm{z}} = \frac{v^2}{r}
        \]
        also hier
        \[
        a_{\mathrm{z}} = \frac{4}{0.5}\,\si{\meter\per\second^2} = \SI{8}{\meter\per\second^2}
        \]
        Richtung: zum Mittelpunkt.
    \end{solutionbox}

\end{enumerate}

\vspace{1cm}

\end{document}
